\documentclass[11pt,a4paper]{article}
\usepackage[margin=23mm]{geometry}
\usepackage{fontspec}\setmainfont{Latin Modern Roman}
\setmonofont{Latin Modern Mono}[Scale=0.85]
\usepackage{amsmath,amssymb,amsthm,longtable,booktabs,xurl,hyperref}
\setlength{\parindent}{0pt}\setlength{\parskip}{6pt}
\setlength{\emergencystretch}{3em}
\newcommand{\code}[1]{\nolinkurl{#1}}
\newcommand{\Adv}{\operatorname{Adv}}
\newcommand{\clip}{\operatorname{clip}}
\newcommand{\PhiB}{\operatorname{Phi}}
\newcommand{\E}{\mathbb E}
\newtheorem{proposition}{Propozycja}
\title{Conditional QROM security of FT1536\\\large Teza docelowa i księga warunków / S05 TARGET\_CONTRACT\_001}
\author{Kontrakt roboczy dla Niirmata}\date{10 października 2026}
\begin{document}\maketitle
\textbf{WARUNKOWY / kontrakt drogi do dowodu bezpieczeństwa.}
Pierwsza część jest samodzielnym kontraktem matematycznym. Stan dowodów,
piny i odwołania do repozytorium są w dodatkach. Warunkowa poprawność
schematu kompozycji nie oznacza, że FT1536 spełnia wszystkie jego warunki.
\textbf{Cel końcowy przyjęty przez właściciela: dowód bezpieczeństwa FT1536
w QROM z deklaracją na takim samym poziomie rygoru jak cel ROM.}
Nie kończymy pracy na mapie, wynikach h=0 ani na niezrealizowanym
warunkowym interfejsie samplera. Bieżący status pozostaje częściowy.

Docelowa deklaracja po domknięciu obowiązków ma brzmieć:
\emph{przy jawnych założeniach kwantowej trudności ISIS/MT-ISIS i prywatnego
generatora oraz w zdefiniowanym modelu wykonania profil FT1536 spełnia
EUF-CMA w QROM, z udowodnioną redukcją i podaną nierównością strat.}
Założenia kryptograficzne pozostają częścią gotowego twierdzenia;
brakujące lematy o tym konkretnym schemacie muszą zostać udowodnione,
a nie zachowane jako zastępcza deklaracja bezpieczeństwa.

\section{Eksperyment i dziedzina twierdzenia}
Badamy EUF-CMA z kwantowymi zapytaniami do publicznego losowego hasza
oraz \emph{klasycznymi} zapytaniami Sign. To standardowy model QROM dla
podpisów; dostęp do Sign w superpozycji jest osobnym, silniejszym celem,
poza tą tezą. Wersja modelu odpowiada konwencjom KLS i Unruha wskazanym
w dodatku bibliograficznym.

Jeden klucz publiczny $h$ jest losowany raz z jawnego prawa $\mu_H$.
W instancji schematu ma to być marginal prawa klucza wyemitowanego przez
rzeczywisty KeyGen po właściwym warunkowaniu na akceptację.
Prawo warunkowe musi być dobrze określone: dodatnia masa akceptacji/emisji
jest jawnym warunkiem instancji, nie domniemaniem z samej nazwy KeyGen.
Przeciwnik otrzymuje klucz publiczny, bez sekretu, seedu i liczby prób KeyGen.
Nie losujemy nowego klucza przy Sign i nie warunkujemy klucza ponownie
na sukcesie podpisania. To samo $\mu_H$ występuje w założeniu trudności.

Budżet przeciwnika to $\beta=(q_s,q_H,t,w,L)$: liczba klasycznych żądań
podpisu, liczba kwantowych zapytań hasza, czas, pamięć i limit długości
wejść. Liczniki reduktora obejmują dodatkowo jego własne zapytania,
implementację orakulum, prywatną losowość, prekomputację i pamięć.
Nie wolno przyznać mu darmowego dodatkowego orakulum spoza zadania trudnego.

Profil współczynnikowy ma $q=18433$, $n=1536$, $\sigma=768$,
$B=2093922385$ i pierścień $R_q=(\mathbb Z/q\mathbb Z)[X]/(X^{1536}-X^{768}+1)$.
Dla wektora $v$ piszemy
\[
 Q_0(v)=\sum_{i=0}^{767}(v_i^2+v_iv_{i+768}+v_{i+768}^2),\quad
 Q(u,v)=Q_0(u)+Q_0(v),\quad A_h(u,v)=u+hv\pmod q.
\]
Hasz $H:\mathcal X\to R_q$ jest jedną wspólną, jednostajnie losową funkcją
na skończonej dziedzinie wejść określonej limitem $L$.
Kwantowe zapytanie ma działanie
$|x,y\rangle\mapsto|x,y\oplus\mathrm{enc}(H(x))\rangle$ przy ustalonym
kodowaniu wyjścia; konwencja kontrolowanych zapytań i koszt jej symulacji
są jawne. Jednostajnego $R_q$ nie utożsamiamy z jednostajnymi bitami kodowania.

\subsection*{Rdzeń współczynnikowy Q-COEFF}
Sign zapisuje całą wiadomość $m$ w SeenSign także przy późniejszej porażce,
losuje jeden idealny nonce $r\in\{0,1\}^{320}$ i ustala $c=H(r\Vert m)$.
Podczas jednego Sign nie ma interaktywnego ujawnienia nonce ani zapytań
przeciwnika pomiędzy jego wewnętrznymi krokami.

Niech $\mathcal V=[-65535,65535]^{1536}\cap\mathbb Z^{1536}$.
Jedna próba pobiera $(u,v)\in\mathcal V^2$ z włókna $A_h(u,v)=c$
z wagą proporcjonalną do $\exp(-Q(u,v)/(2\sigma^2))$.
Gdy włókno jest puste, próba daje porażkę. Próbę odrzucamy, jeśli
$Q(u,v)\ge B$. Dopuszczamy najwyżej 16 niezależnych prób przy tym samym $c$.
Po pierwszej akceptacji następuje jedna emisja: zwracamy $v$, jeśli
wszystkie współczynniki są w $[-32768,32767]$, a w przeciwnym razie porażkę.
Brak akceptacji po 16 próbach również daje porażkę. Cała odpowiedź to
$(r,o)$ z $o$ będącym podpisem albo znacznikiem porażki $\bot$.
Nie warunkujemy prawa na wydaniu podpisu i nie poprawiamy go dodatkowym Verify.

Weryfikacja sprawdza signed16 i $Q(\mathrm{center}_q(c-hs),s)<B$,
z $c=H(r\Vert m)$. Wygrana to zaakceptowane klasyczne fałszerstwo
$(m^*,r^*,s^*)$, gdzie $m^*\notin\mathrm{SeenSign}$.
Nie twierdzimy tu silnej niepodrabialności sEUF-CMA.

\subsection*{Wersja programu Q-API}
W Q-API prywatne KeyGen/Sign/Verify, serializacje i generator są rzeczywiste
w przypiętym profilu wykonania. Publiczne H2P pozostaje zastąpione
wspólnym orakulum QROM; konkretna publiczna funkcja SHAKE i jej H2P
wymagają osobnego rozszerzenia. Obserwowane odpowiedzi rozróżniają
PRE\_ABORT, $(r,\mathrm{POST\_ABORT})$ oraz $(r,\mathrm{bytes})$.
Brak powrotu nie jest automatycznie odpowiedzią timeout, a STUCK nie jest
źródłowym abortem. Wersja programu wymaga osobnego powiązania tych zachowań
z Q-COEFF. Funkcyjny model nie obejmuje timing, liczby retry ani innych
kanałów bocznych. Usługi środowiska muszą mieć tę samą określoną semantykę
we wszystkich porównywanych grach.

\section{Zdanie główne: warunkowa redukcja z jawnymi stratami}
Ustalmy zbiór kluczy $K$ zawierający cały nośnik $\mu_H$ i parametry
wyeksportowane przez wymienione poniżej certyfikaty. Obecny mocny cel
samplera ma $K=R_q$; inne wybranie $K$ nie następuje przez sam zapis tego zdania.

\textbf{Jeżeli} zachodzą Q-MODEL oraz certyfikaty Q-KEY, Q-SAMPLER,
Q-HASH-LOG, Q-COMPOSE, Q-COLL, Q-EMBED i Q-RESC,
\textbf{to} dla każdego dopuszczalnego kwantowego $A$ istnieje efektywny
reduktor $B_A$ o zasobach $R_B(\beta)$ i liczbie celów $T_Q(\beta)$ taki, że
\[
 \varepsilon_{\mathrm{coeff}}\le
 \clip\!\left[\Delta_{\mathrm{QRO}}+
 \PhiB\!\left(D,\clip\!\left[
 L_Q\Adv^{QMT\text{-}ISIS}_{\mu_H,T_Q}(B_A)
       +\delta_{\mathrm{emb}}+q_s\eta\right]\right)\right],
 \qquad D=(1+e)^{q_s}-1.                         \tag{TQ-core}
\]
Tutaj $\clip(x)=\min\{1,x\}$ dla nieujemnych $x$, a
\[
 \PhiB(D,b)=\frac{2b+D+\sqrt{D^2+4Db(1-b)}}{2(1+D)}.
\]
W zadaniu QMT-ISIS klucz pochodzi z tego samego $\mu_H$, cele są niezależne
jednostajne w $R_q$, a reduktor ma podać indeks celu i krótki przeciwobraz
spełniający $A_hz=c_j$ oraz $Q(z)<B$.
$T_Q=1$ jest dopuszczalną trasą. $T_Q=q_H+1$ nie jest przesłanką domyślną.

Jeżeli dodatkowo Q-BIND i Q-PRG dostarczają \emph{konkretnych addytywnych}
porównań zdarzenia wygranej pomiędzy Q-API a Q-COEFF, odpowiednio
z granicami $\delta_{\mathrm{bind}}$ i $\delta_{\mathrm{PRG}}^Q$, wówczas
\[
 \varepsilon_{\mathrm{real}}^{\mathrm{QRO}}\le
 \clip\!\left[\delta_{\mathrm{bind}}+\delta_{\mathrm{PRG}}^Q
 +\Delta_{\mathrm{QRO}}+
 \PhiB\!\left(D,\clip\!\left[
 L_Q\Adv^{QMT\text{-}ISIS}_{\mu_H,T_Q}(B_A)
       +\delta_{\mathrm{emb}}+q_s\eta\right]\right)\right]. \tag{TQ-api}
\]
Jeśli source bridge ma postać kierunkowego momentu, mnożnika lub innego
operatora zamiast addytywnego porównania, trzeba zachować jego rzeczywistą
kolejność w kompozycji. TQ-api w powyższej postaci nie jest wtedy dostępne
bez dodatkowego dowodu; nie zakładamy, że wszystkie błędy są addytywne.

$\Delta_{\mathrm{QRO}}$ jest \emph{jednym} udowodnionym kosztem przejścia
od uczciwego orakulum do gry programowanej. Może składać się z kilku
rozłącznych kroków. Nie dopisujemy obok niego starego klasycznego
$\varepsilon_{\mathrm{coll}}$, jeśli ta sama trudność jest już objęta
reprogramowaniem. $L_Q$, $T_Q$, $\delta_{\mathrm{emb}}$ i $R_B$ wynikają
z konstrukcji reduktora, nie z doboru po zaobserwowaniu sukcesu przeciwnika.

W podstawowej trasie $J$ jest prawem wykonalnego S, a $\eta=0$.
W wariancie idealnego $J$ z osobnym przybliżeniem $\widehat J$ wymagamy
pełnego, jednolitego boundu TV $\le\eta$ i wspólnych kanałów kontynuacji;
wtedy $q_s\eta$ leży \emph{wewnątrz} Phi. Nie liczymy tego samego błędu
ponownie w $e$. Kwantowe użycie $D$ wymaga Q-COMPOSE; sama postać potęgi
nie przenosi klasycznego dowodu do QROM.

\section{Księga warunków}
Poniższe warunki mogą być przesłankami twierdzenia warunkowego.
\emph{Ich klasy są różne}: model określa eksperyment, założenia
kryptograficzne dotyczą trudności, a obowiązki dowodowe trzeba dostarczyć
przed ogłoszeniem instancji twierdzenia dla FT1536.
\small
\begin{longtable}{@{}p{.22\linewidth}p{.71\linewidth}@{}}
\toprule Nazwa i klasa & Pełna treść \\\midrule\endhead
Q-MODEL: model & Gra i obserwacje z sekcji 1, kwantowy H i klasyczny Sign,
pojedynczy klucz, jawne budżety, brak ukrytego timing/retry leakage.
Idealne bity i nonce dotyczą tylko jawnej warstwy idealnej.\\
A-QMT: krypto & Dla dokładnego $\mu_H$, liczby celów, normy i zasobów
reduktora zakładamy granicę przewagi kwantowego solvera ISIS/MT-ISIS.
To osobne założenie; koszt samplera nie dowodzi trudności kratowej.\\
Q-PRG: krypto i dowód hopu & Przewaga kwantowych rozróżniaczy prywatnego
rzeczywistego generatora na rzeczywistym horyzoncie. Dokładny sposób dostępu
jest taki jak w grze. Trzeba skonstruować rozróżniacze i policzyć ich zasoby.
Nie przyznajemy dostępu do sekretnego seedu ani darmowej idealizacji klucza.\\
Q-KEY: obowiązek & Identyfikacja $\mu_H$ z prawem zaakceptowanych/emitted
kluczy, dokładne warunkowanie, jeden wspólny klucz w całej grze.
$K$ obejmuje cały potrzebny nośnik; nie usuwa się złych kluczy bez rachunku.\\
Q-SAMPLER: obowiązek & Jeden publiczny, wykonalny S z opisanym kosztem,
bez trapdooru i bez odczytywania ukrytych zapytań kwantowych.
Dla każdego wymaganego h i dozwolonej historii jego pełne prawo
$J_h(c,o)$, łącznie z każdą porażką, spełnia $J_h\ll P_h$ i
$\sum J_h^2/P_h\le1+e$. $e$ ma wynikać z dowodu.
Prywatny stan samplera jest ujęty w certyfikacie lub eliminowany;
wywołanie ze stałym publicznym initial jest dopuszczalną realizacją.\\
Q-HASH-LOG: obowiązek & Legalna kwantowa semantyka wyroczni i stanu reduktora,
spójne ramki i reprogramowanie. Nazwa nie daje uprawnienia do mierzenia
lub klasycznego zapisywania wszystkich zapytań przeciwnika.\\
Q-COMPOSE: obowiązek & Wspólne warunkowe kanały/stan po tych samych
klasycznych losowaniach i adaptacyjny rachunek pełnych momentów.
Musi prowadzić do dokładnego $D=(1+e)^{q_s}-1$ w zadanych grach.\\
Q-COLL: obowiązek & Osobne oszacowanie kwantowej zmiany orakulum,
w tym reprogramowania i ewentualnych kolizji. Entropia nonce i wszystkie
zapytania są liczone we właściwym momencie; bez podwójnego liczenia strat.\\
Q-EMBED: obowiązek & Konstrukcja osadzenia celów i ekstrakcji,
właściwy rozkład celu oraz strata $L_Q$ i $\delta_{\mathrm{emb}}$.
Liczba celów $T_Q$ jest wnioskiem konstrukcji. Samo algebraiczne
wydobycie świadka z zaakceptowanego podpisu nie daje tego kroku.\\
Q-BIND: obowiązek & Powiązanie praw i obserwacji rzeczywistego programu
z grami modelowymi: bajty, Verify, wszystkie aborty, RNG, STUCK i nonreturn.
Jednolitość względem kwantowej kontynuacji musi być wykazana.
To nie założenie „kod jest poprawny”.\\
Q-RESC: obowiązek & Czas, pamięć, zapytania i losowość B, S, prekomputacji,
monet i obsługi orakulum. Skończony, lecz wykładniczy cap nie jest
użytecznym certyfikatem efektywności.\\
S04: osobne rozszerzenie & Publiczny bit-output XOF/H2P i interakcje
z prywatnym użyciem. Twierdzenie o idealnym H nie staje się automatycznie
twierdzeniem o konkretnym publicznym SHAKE.\\
\bottomrule\end{longtable}\normalsize

Nie wolno zamienić otwartego Q-SAMPLER, Q-EMBED lub Q-BIND w małą
liczbę przyjętą bez dowodu. TQ-core i TQ-api opisują warunkową implikację;
pełna instancja wymaga świadków tych warunków, nie tylko listy ich nazw.

\section{Uzasadnienie kolejności strat i kryteria domknięcia}
Schemat dowodu ma następujące krawędzie:
\[
 G_{\mathrm{real}}\ \longrightarrow\ G_{\mathrm{coeff}}
 \ \longrightarrow\ G_{\mathrm{programmed}}
 \ \longrightarrow\ G_{\mathrm{ideal\ sim}}
 \ \longrightarrow\ G_{\mathrm{executed\ sim}}
 \ \longrightarrow\ \mathrm{QMT\text{-}ISIS}.
\]
Pierwsza krawędź konsumuje Q-BIND/Q-PRG w ich wykazanej postaci;
druga Q-HASH-LOG/Q-COLL; trzecia Q-SAMPLER/Q-COMPOSE przez Phi;
czwarta opcjonalne $\eta$; ostatnia Q-EMBED/Q-RESC.
Z monotoniczności Phi w drugim argumencie wynika kolejność TQ-core.
Przewaga jest co najwyżej 1, co uzasadnia końcowe obcięcie.
To algebra składania krawędzi, a nie dowód istnienia brakujących krawędzi.

Etap uznajemy za domknięty dopiero, gdy określono grę i prawo kluczy,
wykonano i związano S, udowodniono jego pełny certyfikat, zbudowano
kwantowy reduktor i zasoby, a wszystkie użyte dowody mają wymagany
status formalizacji i niezależnego odbioru. Twierdzenie o kodzie wymaga
nadto źródłowych bridge'ów. Warunkowy rdzeń nie zastępuje tej instancji.
Nie podajemy tu żadnej liczby bitów bezpieczeństwa. Użyteczność całego
boundu będzie osobno oceniana dla jawnego budżetu zapytań i primitive.

\section{Kierunek dalszej pracy i decyzja o silniejszym celu}
Teza bezpieczeństwa używa określonego prawa $\mu_H$. Certyfikat na całym
$R_q$ jest wystarczający, lecz silniejszy od certyfikatu na całym nośniku
tego prawa. Klucze poza nośnikiem nie występują w eksperymencie.
Samo to spostrzeżenie nie dostarcza charakterystyki nośnika ani nowego
typu formalnego certyfikatu.

\textbf{Główny cel: bezpieczeństwo FT1536 z jego rzeczywistym prawem KeyGen.}
To znaczenie żądanej przez właściciela deklaracji bezpieczeństwa.
Jednolity all-h sampler zachowujemy jako osobny, mocniejszy cel pomocniczy,
który nadal ma status OPEN. Jego historyczny kwantyfikator nie został
zmieniony i nie ogłaszamy go rozwiązanego przez certyfikat na mniejszym zbiorze.
Do głównego twierdzenia wolno użyć odpowiednio udowodnionego certyfikatu
na całym nośniku rzeczywistego prawa kluczy; charakterystykę nośnika
i nowe interfejsy trzeba dopiero wykazać. Nie wykluczamy h=0 lub h=1
przez sam fakt, że stanowią trudny przykład dla konkretnego samplera.

Najbliższy obowiązek to Q-SAMPLER wraz z Q-RESC: adaptowana propozycja
lub kontrolowana faktoryzacja, z nowym certyfikatem pełnej odpowiedzi.
Wysoka precyzja normalizatora nie naprawia złego rozkładu ani akceptacji
wykładniczo małej. Równoległym logicznie, ale osobno zlecanym podcelem
Q-EMBED jest przypadek $q_s=0$; nie utożsamiamy go z pełnym EUF-CMA.
Kernelizacja wybranych lematów jest osobnym obowiązkiem, nie substytutem
brakującej instancji samplera.

\appendix
\section{Why simpler samplers fail — dokładny zakres przeszkód}
\begin{proposition}[DyadicObstruction — tekstowy]
Deterministyczna funkcja co najwyżej b uczciwych bitów nie daje dokładnie
jednostajnego współczynnika modulo 18433. Każda masa ma postać $k/2^b$,
a równość z $1/18433$ wymagałaby $18433\mid2^b$.
\end{proposition}
To przeszkoda dla dokładnej jednostajności przy ustalonej taśmie;
nie wyklucza aproksymacji ani losowania o nieograniczonej długości.
Szkic Lean nie jest potwierdzonym wynikiem kernelowym.

\begin{proposition}[Niepoprawny marginal sekwencyjny — tekstowy z balls]
Dla opisanej w NORMALIZER-005 konstrukcji i h=1,c=0 moment warunkowy
\emph{przed normą} przekracza $2^{1536}$.
\end{proposition}
Nie jest to dolna granica końcowego momentu po normie, cap16, emisji
lub uśrednieniu wyzwania. Wynik dotyczy tej konstrukcji.

\begin{proposition}[Odrzucanie całych wektorów — tekstowy z balls]
Dla idealnej propozycji U=w/G, stałego envelope i h=1,c=0 w MARGINAL-006
akceptacja jest mniejsza niż $2^{-767}$. Sukces w $2^{128}$ propozycjach
ma prawdopodobieństwo mniejsze niż $2^{-639}$.
\end{proposition}
To przeszkoda kosztowa tej strategii, nie dowód niemożliwości faktoryzacji,
innych propozycji ani QROM dla FT1536. Liczby opisują rachunek kosztu,
nie poziom bezpieczeństwa schematu.

\section{Stan dowodów i mapowanie na repozytorium}
\small
\begin{longtable}{@{}p{.30\linewidth}p{.63\linewidth}@{}}
\toprule Wynik & Status i rola \\\midrule\endhead
Algebra ekstrakcji & Istniejący eksport Lean: Relation.accepted\_extracts.
Nie zastępuje Q-EMBED. Nowego replayu tego modułu tu nie wykonano.\\
CONTINUATION-002 & Tekstowa kompozycja i model reprogramowania.
Instancja kwantowa, formalizacja i zasoby pozostają otwarte.\\
SAMPLER-004, h=0 & Tekstowy pełny moment z rachunkiem Sage.
Brak kernela; nie jest to wynik dla całego emitted key law.\\
NORMALIZER-005 & Tekstowy wynik w dokładnym zapisanym zakresie.
Wiążące liczby pochodzą z CORRECTED certyfikatu i erratum końców MPFR.\\
MARGINAL-006 & Nowy warunkowy certyfikat pełnej mieszanki capu;
nie dostarcza efektywnej instancji samplera dla wszystkich h.\\
DyadicObstruction & Tekstowy dowód i nieuruchomiony szkic Lean.
Brak potwierdzenia kernela; wcześniejszy timeout dotyczył innego modułu.\\
Q-SAMPLER/Q-EMBED/Q-BIND/Q-RESC & Otwarte w pełnym wymaganym zakresie;
nie ma kompletnego twierdzenia QROM dla rzeczywistego programu.\\
\bottomrule\end{longtable}\normalsize

Piny wejść: \code{TARGET_CONTRACT_001_INPUTS.json} i
\code{REVIEW_006_INPUTS.json}. Zakres bieżącego przeglądu drugiego okna,
replay i ograniczenia: \code{REVIEW_006_001.tex} oraz jego receipt.
Nie nadano statusu REVIEWED i nie zmieniono T12.1/B20.

Interfejsy źródłowe: \code{Run2/LawBinding.lean},
\code{Run2/ConcreteReduction.lean}, \code{FT1536/Relation.lean};
wcześniejszy kontrakt programu: \code{T12_1/END_TO_END_SCOPE.md},
\code{FT1536_M0_CONTRACT_RUN_001/GAME.md} i \code{TARGET_TYPE.md}.
Mocny \code{LocalJointCertificate} kwantyfikuje po całym $R_q$;
nie zmieniono go w tym pakiecie.

\section{Źródła modelu QROM}
Kiltz--Lyubashevsky--Schaffner, zaakceptowany manuskrypt 20 II 2018:
\url{https://pure.uva.nl/ws/files/45974858/916.pdf}.
Unruh, wersja konferencyjna EUROCRYPT 2015:
\url{https://www.iacr.org/archive/eurocrypt2015/90560105/90560105.pdf}.
Wersje i wcześniejszy zakres ekstrakcji zapisano w CONTINUATION-002.
Ten kontrakt nie ogłasza nowości ani niezależnego audytu całych publikacji.
\end{document}
